\documentclass[10pt]{extarticle}
\usepackage[a4paper,margin=1.5cm]{geometry}
\usepackage[T1]{fontenc}
\usepackage{amsmath,amssymb}
\usepackage{graphicx}
\usepackage[export]{adjustbox}
\usepackage{multicol}
\usepackage{enumitem}
\usepackage{booktabs,array,tabularx}
\setlength{\columnsep}{0.7cm}
\setlength{\parindent}{0pt}
\setlength{\parskip}{3pt}
\pagestyle{empty}
\makeatletter
% figure registry (figdefs.tex is generated by the extract step)
\newcommand{\qfigdef}[4]{\@namedef{qf@file@#1}{#2}\@namedef{qf@w@#1}{#3}\@namedef{qf@wide@#1}{#4}}
\newcommand{\qf@check}[1]{\@ifundefined{qf@file@#1}{\PackageError{score2}{Unknown figure id #1}{}}{}}
\newcommand{\qf@img}[4][]{\includegraphics[width=#2\dimexpr\@nameuse{qf@w@#4}pt\relax,max width=#3,max height=0.45\textheight,#1]{\@nameuse{qf@file@#4}}}
% Full-width questions: finalize wraps every question that contains a wide figure in \qfullbegin...\qfullend,
% so the WHOLE question (stem, figures, options) is set across the page in reading order, then the two columns
% resume. The question's head - its number, stem and first figure (\qfullsplit marks the end of the head) -
% is kept on one page, so a number is never left at a page foot with its figure overleaf; the rest may break
% across pages normally, so no big gaps are left. Without a \qfullsplit, the whole question is kept together
% if it is under 45% of a page. Never write these yourself.
\newif\ifqf@full
\newif\ifqf@boxed
\newsavebox{\qf@box}
\newcommand{\qf@open}{\global\qf@fulltrue\par\global\qf@boxedtrue
  \setbox\qf@box\vbox\bgroup\hsize\textwidth\linewidth\textwidth}
\newcommand{\qf@emit}[1]{\par\egroup\global\qf@boxedfalse
  \ifdim\dimexpr\ht\qf@box+\dp\qf@box\relax<#1\textheight\noindent\usebox{\qf@box}\par\else\unvbox\qf@box\fi}
\newcommand{\qfullsplit}{\ifqf@boxed\qf@emit{0.85}\fi}
\newcommand{\qf@close}{\ifqf@boxed\qf@emit{0.45}\fi\global\qf@fullfalse}
\newcommand{\qfullbegin}{\end{multicols}\qf@open}
\newcommand{\qfullend}{\qf@close\begin{multicols}{2}}
\newcommand{\qfullstart}{\qf@open}
\newcommand{\qfulldone}{\qf@close}
% \qfig[scale]{id}: block figure, centred, directly where it is written (size set by the tool).
% Fallback if a wide figure is not inside a full-width question: break out of the columns in place - no floats.
\newcommand{\qfig}[2][1]{\qf@check{#2}%
  \ifqf@full\qf@blk{#1}{#2}\else
  \ifnum\@nameuse{qf@wide@#2}=1 \expandafter\@firstoftwo\else\expandafter\@secondoftwo\fi
  {\end{multicols}\begin{center}\qf@img{#1}{\textwidth}{#2}\end{center}\begin{multicols}{2}}%
  {\qf@blk{#1}{#2}}\fi}
\newcommand{\qf@blk}[2]{\par\smallskip{\centering\qf@img{#1}{\linewidth}{#2}\par}\smallskip}
% \qfigi[scale]{id}: inline figure (inside option lists / tables / minipages); never breaks columns.
\newcommand{\qfigi}[2][1]{\qf@check{#2}\qf@img[valign=c]{#1}{\linewidth}{#2}}
\makeatother
% structure macros - lint reads these to check completeness and figure placement
\newcommand{\papersec}[1]{\par\bigskip{\centering\Large\bfseries PAPER-#1\par}\medskip}
\newcommand{\partsec}[2]{\par\medskip\noindent{\large\bfseries PART-#1: #2}\par\smallskip}
\newcommand{\sect}[1]{\par\smallskip\noindent\textbf{#1}\par}
\newcommand{\qn}[1]{\par\smallskip\noindent\textbf{#1.}~}
\newcommand{\paraq}[2]{\par\smallskip\noindent\textbf{Paragraph for Questions #1 and #2}\par}
\newenvironment{qoptions}{\begin{itemize}[leftmargin=1.8em,itemsep=1pt,topsep=2pt,label={}]}{\end{itemize}}
\newcommand{\opt}[2]{\item[(#1)] #2}
% figure registry (generated by score2_tools.py finalize)
\qfigdef{p1_phy_s0_q01_1}{SRGBT - Test 10 (26 Sep 2021)-img/p1_phy_s0_q01_1.png}{87.2}{0}
\qfigdef{p1_phy_s0_q01_4}{SRGBT - Test 10 (26 Sep 2021)-img/p1_phy_s0_q01_4.png}{136.4}{0}
\qfigdef{p1_phy_s0_q01_3}{SRGBT - Test 10 (26 Sep 2021)-img/p1_phy_s0_q01_3.png}{112.7}{0}
\qfigdef{p1_phy_s0_q04_1}{SRGBT - Test 10 (26 Sep 2021)-img/p1_phy_s0_q04_1.png}{114.0}{0}
\qfigdef{p1_phy_s0_q05_1}{SRGBT - Test 10 (26 Sep 2021)-img/p1_phy_s0_q05_1.png}{152.3}{0}
\qfigdef{p1_phy_s0_q09_1}{SRGBT - Test 10 (26 Sep 2021)-img/p1_phy_s0_q09_1.png}{86.5}{0}
\qfigdef{p1_phy_s0_q11_1}{SRGBT - Test 10 (26 Sep 2021)-img/p1_phy_s0_q11_1.png}{72.5}{0}
\qfigdef{p1_phy_s0_q12_1}{SRGBT - Test 10 (26 Sep 2021)-img/p1_phy_s0_q12_1.png}{128.1}{0}
\qfigdef{p1_phy_s0_q13_1}{SRGBT - Test 10 (26 Sep 2021)-img/p1_phy_s0_q13_1.png}{81.7}{0}
\qfigdef{p1_phy_s0_q14_1}{SRGBT - Test 10 (26 Sep 2021)-img/p1_phy_s0_q14_1.png}{110.6}{0}
\qfigdef{p1_phy_s0_q15_1}{SRGBT - Test 10 (26 Sep 2021)-img/p1_phy_s0_q15_1.png}{112.6}{0}
\qfigdef{p1_phy_s0_q19_1}{SRGBT - Test 10 (26 Sep 2021)-img/p1_phy_s0_q19_1.png}{73.7}{0}
\qfigdef{p1_phy_s0_q20_1}{SRGBT - Test 10 (26 Sep 2021)-img/p1_phy_s0_q20_1.png}{64.9}{0}
\qfigdef{p1_chem_s0_q23_1}{SRGBT - Test 10 (26 Sep 2021)-img/p1_chem_s0_q23_1.png}{164.4}{0}
\qfigdef{p1_chem_s0_q31_1}{SRGBT - Test 10 (26 Sep 2021)-img/p1_chem_s0_q31_1.png}{51.3}{0}
\qfigdef{p1_chem_s0_q39_1}{SRGBT - Test 10 (26 Sep 2021)-img/p1_chem_s0_q39_1.png}{68.3}{0}
\raggedcolumns
\makeatletter\@beginparpenalty=10000\makeatother
\newcolumntype{L}[1]{>{\raggedright\arraybackslash}p{#1}}
\begin{document}
\begin{center}{\Large\bfseries SRGBT --- Test 10 (26 Sep 2021)}\end{center}
\begin{multicols}{2}
\raggedcolumns
\papersec{1}
\partsec{1}{PHYSICS}
\qn{1} To one end of a light inextensible string of length $\ell$ is attached a particle of mass m resting on a smooth horizontal a particle of mass m resting on a smooth horizontal table. The string passes through a smooth hole in the table and to its other end is attached a small particle of same mass. The system is set in rotation as shown in figure with constant angular speed. Choose the correct option(s):\par
\qfig{p1_phy_s0_q01_1}
\begin{qoptions}
\opt{A}{$\dfrac{\ell_1}{\ell_2}=\dfrac{\omega_2^2}{\omega_1^2}$}
\opt{B}{$\dfrac{\ell_1}{\ell_2}=\dfrac{\omega_1^2}{\omega_2^2}$}
\opt{C}{The motion is possible if, $\dfrac1{\omega_1^2}+\dfrac1{\omega_2^2}<\dfrac{\ell}{g}$}
\opt{D}{The motion is possible if, $\dfrac1{\omega_1^2}+\dfrac1{\omega_2^2}>\dfrac{\ell}{g}$}
\end{qoptions}
\paraq{2}{3}Wind velocities at various height can be investigated using balloons with constant ascending velocity. A graph between angle of elevation ($\theta$) against time is shown in the diagram by observing such a balloon. The balloon was released at distance 900 m from the point of observation from the same level and it seemed to be rising directly upwards. Wind velocity near the ground was zero. Wind velocity near the ground was zero. Wind velocity is horizontal at each level.\par
\qfig{p1_phy_s0_q01_4}
\qfig{p1_phy_s0_q01_3}
\qn{2} Choose the correct statement(s) :
\begin{qoptions}
\opt{A}{During the journey of the balloon its distance from point of observation always decreases.}
\opt{B}{During the journey of the balloon its distance from point of observation increases.}
\opt{C}{Angular velocity of balloon with respect to point of observation is 3.7 deg/sec. at t = 10 sec.}
\opt{D}{Angular velocity of balloon with respect to point of observation is zero at t = 10 sec.}
\end{qoptions}
\qn{3} If slope of tangent of graph at t = 0, i.e., $\tan\phi=\dfrac6\pi$ deg/sec., then :
\begin{qoptions}
\opt{A}{Angular velocity of balloon about point of observation is zero at t = 0.}
\opt{B}{Velocity of ascending of balloon is 30 m/s}
\opt{C}{Velocity of wind at height 300 m is 40 m/s}
\opt{D}{Height of balloon at t = 10 sec is 40 m.}
\end{qoptions}
\qn{4} A block of mass m is released on a track shown in figure. The track is smooth from A to D and rough from point D onwards. Based on above information, answer the following questions.\par
\qfig{p1_phy_s0_q04_1}
\begin{qoptions}
\opt{A}{The normal force N$_\mathrm B$ exerted by the track just after crossing point B is 4 mg}
\opt{B}{The normal force N$_\mathrm C$ at C is 7 mg}
\opt{C}{The normal force N$_\mathrm B$ exerted by the track just after crossing point B is 7 mg}
\opt{D}{The distance travelled along the incline beyond point D where the block stops is $\dfrac{4R}{\sqrt3\,\mu_\mathrm k+1}$}
\end{qoptions}
\qn{5} A small object of mass m = 1 kg is released from rest at the top of an inclined plane that connects to a horizontal plane without an edge. It slides onto a cart of mass M that has a semi-cylindrical surface of radius R = 0.36 m fixed to the middle of it, as shown in the figure.\par
\qfig{p1_phy_s0_q05_1}
The small object reaches the topmost point of the semi-cylinder and stops there. In continues to move vertically with free fall and hits the cart exactly at the edge. All friction can be neglected.\par
Choose correct option(s) :
\begin{qoptions}
\opt{A}{the minimum possible length of the cart is 1.44 m}
\opt{B}{the mass of the cart is 2.56 kg}
\opt{C}{height h was the small object released is 1.81 m}
\opt{D}{the mass of the cart is 5 kg}
\end{qoptions}
\qn{6} A small object of mass m is suspended by a thread of length $\ell$ and released when thread is horizontal.
\begin{qoptions}
\opt{A}{Velocity of the object when vertical component of its velocity is maximum $\sqrt{\dfrac2{\sqrt3}\,g\ell}$.}
\opt{B}{Tension in the thread when vertical component of velocity is maximum $\sqrt3\,mg$.}
\opt{C}{Ratio of maximum vertical velocity to maximum horizontal velocity is $\sqrt{\dfrac2{3\sqrt3}}$.}
\opt{D}{Ratio of maximum vertical velocity to maximum horizontal velocity is $\sqrt{\dfrac2{\sqrt3}}$.}
\end{qoptions}
\qn{7} A particle is projected in air from origin with speed u and at an angle $\theta$ with positive x-axis. At its maximum height its trajectory has radius of curvature R and centre of curvature has co-ordinate $\left(20\text{m},\dfrac{70}{3}\text{m}\right)$. Choose the correct option(s). ($\vec g=-10\text{m/s}^2\,\hat j$)
\begin{qoptions}
\opt{A}{Equation of trajectory of the projectile motion is $y=3x-\dfrac{3x^2}{40}$.}
\opt{B}{The maximum height of the projectile is 30 m.}
\opt{C}{The value of u is $\sqrt{\dfrac{2000}3}$ m/s.}
\opt{D}{The value of $\theta$ is $\tan^{-1}\left(\dfrac3{\sqrt{10}}\right)$.}
\end{qoptions}
\qn{8} A particle moves along a circular path of radius 3m about a stationary axis according to the law $\theta=(t^3-3t^2)\pi$ radian. Then which of the following is/are correct.
\begin{qoptions}
\opt{A}{The magnitude of mean value of the angular velocity over the time interval between t = 0 and the time when the particle comes to rest, is $2\pi$ rad/sec.}
\opt{B}{The magnitude of average speed over the time interval between t = 0 to t = 3 sec is $8\pi$ m/sec.}
\opt{C}{The distance travelled by the particle between the time interval t = 1 to t = 4 sec, is 66$\pi$ meter.}
\opt{D}{The distance travelled by the particle between the time interval t = 1 to t = 4 sec, is 60$\pi$ meter.}
\end{qoptions}
\qn{9} Figure shows a smooth fixed track with two straight parts connected by a semicircular arc ABC. The shown track is placed in a vertical plane. A small ball is given an initial speed `u' along the track as shown. The ball loses contact with the track somewhere in the part BC and collides with the track at A during subsequent motion. Now answer the following questions.\par
\qfig{p1_phy_s0_q09_1}
\begin{qoptions}
\opt{A}{the value of `u' is $\sqrt{\dfrac72Rg}$}
\opt{B}{time after which it collides at `A' after losing the contact with the shown track, is $\sqrt{\dfrac{6R}g}$}
\opt{C}{the value of `u' is $\sqrt{\dfrac92Rg}$}
\opt{D}{The time after which it collides at `A' after losing the contact with the shown track, is $\sqrt{\dfrac{3R}g}$}
\end{qoptions}
\qn{10} The potential energy function of a particle is given by $\mathrm U(\mathrm r)=\dfrac{\mathrm A}{2\mathrm r^2}-\dfrac{\mathrm B}{3\mathrm r}$, where A and B are constants and r is the radial distance from the centre of the force. Choose the correct option(s) :
\begin{qoptions}
\opt{A}{The equilibrium distance will be $\mathrm r_0=\dfrac{2\mathrm A}{\mathrm B}$}
\opt{B}{The equilibrium distance will be $\mathrm r_0=\dfrac{3\mathrm A}{\mathrm B}$}
\opt{C}{If the total energy of the particle is $\dfrac{\mathrm B^2}{6\mathrm A}$, then its radial velocity will vanish at $\dfrac{\mathrm r_0}3$, where r$_0$ is the equilibrium distance.}
\opt{D}{If the total energy of the particle is $\dfrac{\mathrm B^2}{6\mathrm A}$, then its radial velocity will vanish at $\dfrac{\mathrm r_0}2$, where r$_0$ is the equilibrium distance.}
\end{qoptions}
\qn{11} A racing track ABCDEFA is such that part ABC is circular of radius 400 m, part CD and FA are straight of length 300 m and part EF is circular of radius 100 m. The coefficient of friction between track and car tyres is 0.1. A racing car always moves on this track with uniform speed. The minimum time (in sec.) that a racing car can take to complete one complete lap is \underline{\hspace{1cm}}.\par
\qfig{p1_phy_s0_q11_1}
\qn{12} Consider a pendulum, consisting of a massless string with a mass on the end. The mass is held with the string horizontal, and then released. The mass swings down, and then on its way back up, the string is cut at point P when it makes an angle of $\theta$ with the vertical. What should be the value of $111\times\cot^2\theta$, so that the mass travels largest horizontal distance from P by the time it return to the height it had when the string was cut ?\par
(g = 10 m/s$^2$)\par
\qfig{p1_phy_s0_q12_1}
\qn{13} A wedge of mass M is kept at rest on an horizontal surface, and a block of mass m is kept on the wedge at the height h from the surface. the angle of the wedgeis $\alpha$, see Fig. There is no friction neither between the block and wedge nor between the surface and the wedge. The system is released into a free motion. Find the time (in sec.) needed for the block to reach the surface. (Given M = 9m, h = 400 m, $\tan\alpha=\dfrac1{\sqrt2}$)\par
\qfig{p1_phy_s0_q13_1}
\qn{14} A tangentially attached slope leads to a circular match-box track with radius R = 32 cm set in a vertical plane. The toy car starts from rest at the top of the slope, runs down the slope and detaches from the track at height h = 3/2 R measured from the bottom. Find the $\left(\dfrac{\text{maximum height}}4\right)$ (in cm) reached by it after it reaches the bottom of the track. (Assume that the toy car is point-like, neglect drag and friction).\par
\qfig{p1_phy_s0_q14_1}
\qn{15} As shown in the figure, a ball is released from a certain height which has to perform circular motion on the vertical track of radius 4 m. The track is absent between points A and B. Calculate the height (in cm) from where the ball has to be released so that it will reach at highest point B of the circular track.\par
\qfig{p1_phy_s0_q15_1}
\qn{16} In a park there are three concentric circular running tracks. Radius of 2nd track is double of first and of 3rd track is triple of first. Three runners are running on these tracks with constant speed. When the runner in the first track completes one round, the runner in 2nd has completed half round and the runner in third track has completed quarter round. If the accelerations of the runners are in ratio $\alpha:\beta:\gamma$, where $\alpha$, $\beta$, and $\gamma$ are least integers, then find the value of $4(\alpha+\beta+\gamma)$.
\qn{17} Two identical cubes of mass 1.5M each are kept almost touching the faces on a smooth horizontal surface. On top of them we neatly placed a smooth ball of mass M, which begins to move vertically downwards, pushing the cubes sideways. The initial velocity of the ball is negligible. Radius of the ball is 5m, side of cube 2.5m. If velocity of ball just before impact with ground is V. Find $\dfrac{3\mathrm V^2}4$.
\qn{18} A rod is rotating with an angular velocity of $3\pi t^3$ rad/s. The angle rotated by the rod from t = 0 to the time when it has an angular velocity of $24\,\pi$ rad/s is given by $\alpha\,\pi$ radian. Fill $\dfrac\alpha2$ in OMR sheet.
\qn{19} Two objects A and B of mass m are connected with a uniform rope of length $\ell$ = 0.3 m and mass m as shown in the figure. The rope can slide in a tube which forms a quarter of a circle, and which is fixed to the edge of the table. Initially 1/3 of the rope hangs vertically. Friction is negligible. If the system is released from rest, and B descends a distance $\ell/3$, what is the speed (in m/s) of B?\par
\qfig{p1_phy_s0_q19_1}
\qn{20} A mass m = 1kg, which is free to move on a horizontal frictionless surface, is attached to one end of a massless string that wraps partially around a frictionless vertical pole of radius r (see the top view in figure). You hold on to the other end of the string. At t = 0, the mass has speed v$_0$ in the tangential direction along the dotted circle of radius R shown. Your task is to pull on the string so that the mass keeps moving along the circular path of radius R. If at any instant `t' force applied by hand on string is f$_0$ = 2N then find out centripetal acceleration (in m/s$^2$) of mass. (Given : R = $\dfrac{2\mathrm r}{\sqrt3}$)\par
\qfig{p1_phy_s0_q20_1}
\raggedcolumns
\papersec{1}
\partsec{2}{CHEMISTRY}
\qn{21} Select \textbf{CORRECT} statement(s).
\begin{qoptions}
\opt{A}{O$_3$ and O$^{2-}$ both are diamagnetic}
\opt{B}{Out of O$_2$, O$_2^+$ and O$_2^-$, O$_2$ have maximum magnetic moment (spin only)}
\opt{C}{Out of O$_2$, O$_2^+$ and O$_3$, O$_2^+$ have least O--O bond length}
\opt{D}{Out of O$_2$, O$_2^{2+}$ and O$_2^{2-}$, only O$_2$ is paramagnetic.}
\end{qoptions}
\qn{22} Select the option(s) having correct order of the properties involved.
\begin{qoptions}
\opt{A}{Size in gas phase : Al$^{3+}$ $<$ Ga$^{3+}$ $<$ H$^{-}$}
\opt{B}{IE$_1$ values : Sb $<$ Te}
\opt{C}{At. size : Ag $>$ Hg}
\opt{D}{IE$_1$ values : Bi $>$ Sb}
\end{qoptions}
\qn{23} Consider the boiling point of some molecular hydrides and monoatomic gases\par
\qfig{p1_chem_s0_q23_1}
According to the given information the correct match are :
\begin{qoptions}
\opt{A}{A to D compounds are H$_2$O, H$_2$S, H$_2$Se, H$_2$Te respectively}
\opt{B}{E to H compounds are HF, HCl, HBr, HI respectively}
\opt{C}{I to L compounds are CH$_4$, SiH$_4$, GeH$_4$, SnH$_4$ respectively}
\opt{D}{M to P are Ne, Ar, Kr, Xe respectively}
\end{qoptions}
\qn{24} One mole of `F' atoms are ionised to F$^-$ the energy released is `X' Joules. Then the incorrect statements:
\begin{qoptions}
\opt{A}{X Joules is sufficient to ionise 1 mole of gaseous Cl into Cl$^+$.}
\opt{B}{X Joules is sufficient to ionise 1 mole of gaseous F into F$^+$.}
\opt{C}{X Joules is sufficient to ionise 1 mole of gaseous F into F$^+$ as well as 1 mole of gaseous Cl into Cl$^+$.}
\opt{D}{Less than 1 mole of gaseous F or Cl will be ionised to F$^+$ or Cl$^+$.}
\end{qoptions}
\qn{25} Find correct match.\par
\begin{tabular}{@{}l@{\,}p{1.6cm}@{\,}p{2.3cm}@{\,}p{2cm}@{}}
 & Molecule & Nature of bridge bond & Planar/Non--planar\\
(A) & Al$_2$Cl$_6$ & 3C--4e$^-$ & NP\\
(B) & I$_2$Cl$_6$ & 3C--4e$^-$ & P\\
(C) & Al$_2$(CH$_3$)$_6$ & 3C--2e$^-$ & NP\\
(D) & (BeCl$_2$)$_\mathrm n$ & 3C--4e$^-$ & NP\\
\end{tabular}\par
\qn{26} X, Y and Z are three consecutive elements of second period whose I.P. data (in eV) are available. By using these data specify incorrect statements regarding X, Y and Z.\par
\begin{tabular}{|c|c|c|c|c|}\hline
Element & IP$_1$ & IP$_2$ & IP$_3$ & IP$_4$\\\hline
X & 5.39 & 75.6 & 122.4 & Not possible\\\hline
Y & 8.2 & 25.1 & 37.9 & 259.2\\\hline
Z & 9.3 & 18.2 & 217.6 & 363.7\\\hline
\end{tabular}\par
\begin{qoptions}
\opt{A}{X$_{(g)}$ + e$^-$ $\longrightarrow$ A$^-_{(g)}$ \quad $\Delta$H = positive}
\opt{B}{Z$_\mathrm Y$ $>$ Z$_\mathrm Z$ (Z is atomic number)}
\opt{C}{Oxide of `Y' is amphoteric}
\opt{D}{X$^{+1}$ is stable than X$^{2+}$}
\end{qoptions}
\qn{27} Which of the following statements are correct ?
\begin{qoptions}
\opt{A}{Vapour pressure of H$_2$O is more than H$_2$Te at given temperature.}
\opt{B}{In both N$_2$O$_5$ and N$_2$O$_4$ all N--O bond lengths are equivalent.}
\opt{C}{In both crystalline NaHCO$_3$ and KHCO$_3$, HCO$_3^-$ forms different type of association of anion through hydrogen bond.}
\opt{D}{Amongst B$_2$, C$_2$, N$_2^-$ and O$_2$, N$_2^-$ and O$_2$ on further ionization (losing single electron) forms thermodynamically more stable species.}
\end{qoptions}
\qn{28} Select the option containing correct order of the parameters given
\begin{qoptions}
\opt{A}{Increasing order of acidic strength : Na$_2$O, MgO, CO$_2$, SO$_3$}
\opt{B}{Decreasing order of ionic radii : K$^+$, Na$^+$, Li$^+$, Mg$^{2+}$}
\opt{C}{Decreasing order of boiling points : H$_2$O, HF, NH$_3$}
\opt{D}{Decreasing order of boiling points : He, Ne, Ar, Xe}
\end{qoptions}
\qn{29} X and Y are second period elements. They form two compounds XF$_3$ and YH$_3$ respectively. XF$_3$ is resonance stabilized while YH$_3$ is not. Select the incorrect statement.
\begin{qoptions}
\opt{A}{X--F bond distance in XF$_3$ is smaller than that in XF$_4^-$.}
\opt{B}{X--F bond in XF$_3$ is larger than that in XF$_3^-$ -- YH$_3^+$.}
\opt{C}{XF$_3$ does not have back bonding.}
\opt{D}{YH$_3$ is a Lewis acid.}
\end{qoptions}
\qn{30} Which of the following statements are \textbf{correct} :
\begin{qoptions}
\opt{A}{IE$_2$ of Cu $>$ IE$_2$ of Zn where as IE$_1$ of Cu $<$ IE$_1$ of Zn.}
\opt{B}{$|\Delta$H$_\mathrm{EG}|$ of S $>$ $|\Delta$H$_\mathrm{EG}|$ of O.}
\opt{C}{Reducing nature of Li in aqueous solution is more than that of Na.}
\opt{D}{There is a substantial contraction in size observed on moving from 3$^\mathrm{rd}$ group to 4th group in 6th period of periodic table.}
\end{qoptions}
\qn{31} Consider following molecular species :\par
XeF$_4$, ClO$_4^-$, XeO$_3$, IF$_4^+$, ClF$_3$, SF$_6$, SF$_4$, PCl$_6^-$, XeO$_6^{4-}$, S$_3$O$_6^{2-}$\par
and calculate $\dfrac{(\mathrm x+\mathrm y)}{\mathrm z}$.\par
x = Total number of molecular species in which all bond lengths are not equal\par
y = Total number of species in which `d' orbital which having zero nodal plane participate in hybridisation\par
z = total number of planar species
\qn{32} How many comparison(s) are \textbf{INCORRECT} among given ?\par
(i) H $>$ Li (ionisation energy)\quad (ii) Li $>$ Be (size)\quad (iii) Na $>$ Rb (size)\quad (iv) O $>$ N (size)\quad (v) S $>$ P (ionisation energy)\quad (vi) C $>$ O (size)\quad (vii) B $>$ Al (size)\quad (viii) F $>$ Cl (ionisation energy)\quad (ix) N $>$ O (ionisation energy)
\qn{33} How many of the following has underlined atom in sp$^2$ hybridisation ?\par
\underline{Si}O$_2$, \underline{N}(SiH$_3$)$_3$, \underline{B}F$_3$, \underline{Al}Cl$_3$.6H$_2$O, H$_3$\underline{O}$^\oplus$, B$_3$\underline{N}$_3$H$_6$, \underline{S}O$_2$, \underline{N}O$_2$
\qn{34} Assuming screening by K shell electrons to be same in Li \& Be. Calculate I.E.$_1$ (in eV) of Li if IE$_2$ of Be is 13.6$\times$1.44 eV/atom.
\qn{35} Calculate the bond length of C--X bond, if C--C bond length is 1.54 \AA, X--X bond length is 1.00 \AA\ and electronegativity values of C and X are 2.0 and 3.0 respectively
\qn{36} From the following data calculate \% ionic character in the bond A-B.\par
[1 eV/atom = 96 KJ/mole]\par
$\Delta$H$_{\mathrm{EG}_1}$ of A$^{+1}$ = $-$960 KJ/mole \hfill IE$_1$ of A$^-$ = $-$1.6 eV/atom\par
$\Delta$H$_{\mathrm{EG}_1}$ of B = $-$192 KJ/mole \hfill IE of B = 3.6 eV/atom
\qn{37} (i) Number of $\pi$--bonds in C$_2$ molecule (according to MOT) = x.\par
(ii) Number of ABMO electrons in O$_2^-$ ion = y\par
(iii) Number of ABMO having energy less than $\sigma$2p$_\mathrm z$ in F$_2$ molecule = z.\par
\textbf{Specify :} $\dfrac{2\mathrm x+4\mathrm y-\mathrm z}{600}=\,?$
\qn{38} The number of correct trends is/are :\par
(i) AgF $>$ AgCl $>$ AgBr $>$ AgI (Solubility in water)\par
(ii) TiO$_2$ $<$ SiO$_2$ (Polarising power)\par
(iii) O$_2^-$ $>$ O$_2^{2-}$ (Ionic radii)\par
(iv) Cu $>$ Ag $<$ Au (Ionisation energy)\par
(v) N $<$ O $<$ F $<$ Ne (Z$_\mathrm{eff}$)\par
(vi) La$^{3+}$ $>$ Nd$^{3+}$ $>$ Gd$^{3+}$ $>$ Ho$^{3+}$ $>$ Lu$^{3+}$ (Ionic radii)
\qn{39} Consider the following compounds in liquid form \qfigi{p1_chem_s0_q31_1} SO$_2$, SiCl$_4$, S$_2$, PCl$_2$F$_3$, O$_3$, CrO$_5$, C$_{12}$O$_9$, Cl$_2$.\par
When a charged comb is brought near their following steam how many of them deflection as per the following figure ?\par
\qfig{p1_chem_s0_q39_1}
\qn{40} Lattice energy of how many ionic compounds are greater than NaCl.\par
NaF, KCl, MgO, MgF$_2$, A$\ell$N, CsBr, NaI, Na$_2$O
\raggedcolumns
\papersec{1}
\partsec{3}{MATHEMATICS}
\qn{41} Let a, b, c be the sides of a triangle such that $\dfrac{bc}{b+c}+\dfrac{ca}{c+a}+\dfrac{ab}{a+b}=s$ where s is semi -- perimeter, then the triangle is
\begin{qoptions}
\opt{A}{equilateral}
\opt{B}{isosceles but not right angled}
\opt{C}{right angled but not isosceles}
\opt{D}{obtuse angled}
\end{qoptions}
\qn{42} If $3\sin\alpha\sin\beta+4\cos\beta+5\sin\beta\cos\alpha=5\sqrt2$, where $\alpha,\beta\in\mathrm R$, then
\begin{qoptions}
\opt{A}{$\tan\alpha=\dfrac35$}
\opt{B}{$\tan\alpha=\dfrac53$}
\opt{C}{$\cos\alpha=\dfrac3{\sqrt{34}}$}
\opt{D}{$\tan\beta=\dfrac{\sqrt{34}}4$}
\end{qoptions}
\qn{43} For a $\triangle$ABC with integer side lengths and $\cos\mathrm A=\dfrac{11}{16}$, $\cos\mathrm B=\dfrac78$, $\cos\mathrm C=-\dfrac14$, which of the following statement(s) is/are true?
\begin{qoptions}
\opt{A}{least possible perimeter of triangle ABC is 27}
\opt{B}{least possible perimeter of triangle ABC is 9}
\opt{C}{least area of the triangle ABC is $\dfrac{3\sqrt{15}}4$}
\opt{D}{None of these}
\end{qoptions}
\qn{44} If $\tan^2x+\tan^2y+\cot^2(x+y)=1$, $x,y\in\left(0,\dfrac\pi4\right)$, then :
\begin{qoptions}
\opt{A}{x can only be $\dfrac\pi6$}
\opt{B}{y can only be $\dfrac\pi6$}
\opt{C}{x can take more than 1 value}
\opt{D}{none of these}
\end{qoptions}
\qn{45} In $\triangle$ABC, P is any point inside a triangle such that area of $\triangle$BPC, $\triangle$APC, $\triangle$APB are equal. Line AP cut BC at M, area of $\triangle$PMC is 5 sq. unit then area of $\triangle$ABC can be :
\begin{qoptions}
\opt{A}{20 sq. units}
\opt{B}{25 sq. units}
\opt{C}{30 sq. units}
\opt{D}{10 sq. units}
\end{qoptions}
\qn{46} In $\triangle$ABC \ $\sin(\mathrm A-\mathrm B)=\dfrac35$, $\sin\mathrm A=2\sin\mathrm B$ \ then \ $\dfrac{\mathrm a+\mathrm b}{\mathrm c}=$
\begin{qoptions}
\opt{A}{$\dfrac4{\sqrt5}$}
\opt{B}{$\sqrt{\dfrac25}$}
\opt{C}{$\dfrac3{\sqrt5}$}
\opt{D}{$\dfrac2{\sqrt5}$}
\end{qoptions}
\qn{47} In $\triangle$ABC, if a = 2, B = 60$^\circ$, C = 75$^\circ$, then
\begin{qoptions}
\opt{A}{$\mathrm b=\sqrt6$ units}
\opt{B}{area of $\triangle$ABC $=\left(\dfrac{3+\sqrt3}2\right)$ sq. units.}
\opt{C}{Length of angle bisector of $\angle$B is 2 units}
\opt{D}{Circumradius of the triangle ABC is $\sqrt2$ units}
\end{qoptions}
\qn{48} $\sqrt{\cos2x}+\sqrt{1+\sin2x}=2\sqrt{\sin x+\cos x}$ \ if
\begin{qoptions}
\opt{A}{$\sin x+\cos x=0$}
\opt{B}{$x=2n\pi$}
\opt{C}{$x=n\pi-\dfrac\pi4$}
\opt{D}{$x=2n\pi\pm\cos^{-1}\left(\dfrac{-1}5\right)$}
\end{qoptions}
\qn{49} If in a $\triangle$ABC, a, b, c are in A.P. then it is necessary that :
\begin{qoptions}
\opt{A}{$\dfrac23<\dfrac bc<2$}
\opt{B}{$\dfrac13<\dfrac bc<\dfrac23$}
\opt{C}{$\dfrac23<\dfrac ba<2$}
\opt{D}{$\dfrac13<\dfrac ba<\dfrac23$}
\end{qoptions}
\qn{50} The solution set of $|\sin x|\le|\cos2x|$ contains
\begin{qoptions}
\opt{A}{$\displaystyle\bigcup_{n\in\mathrm I}\left\{\left[n\pi-\dfrac\pi6,n\pi+\dfrac\pi6\right]\right\}$}
\opt{B}{$\displaystyle\bigcup_{n\in\mathrm I}\left\{n\pi+\dfrac\pi2\right\}$}
\opt{C}{$\displaystyle\bigcup_{n\in\mathrm I}\left\{\left[n\pi-\dfrac\pi8,n\pi+\dfrac\pi8\right]\right\}$}
\opt{D}{$\displaystyle\bigcup_{n\in\mathrm I}\left\{\left[n\pi-\dfrac\pi4,n\pi+\dfrac\pi4\right]\right\}$}
\end{qoptions}
\qn{51} A triangle ABC has A $-$ B = 120$^\circ$ and the circumradius is 8 times the in--radius. If $\cos\mathrm C=\mathrm a/8$, then a =
\qn{52} The number of ordered pair (x,y) satisfying equation $\sin x\cos y=1$ and $x^2+y^2\le9\pi^2$ is m the find $\dfrac{78}{2\mathrm m}$
\qn{53} Total number of ordered pairs (x, y) satisfying $x\left(\sin^2x+\dfrac1{x^2}\right)=2\sin x.\sin^2y$, where $x\in(-\pi,0)\cup(0,\pi)$ and $y\in[0,2\pi]$ is / are
\qn{54} The least degree of a polynomial with integer coefficient whose one of roots may be $\cos12^\circ$ is :
\qn{55} If $x=\sin\dfrac{2\pi}7+\sin\dfrac{4\pi}7+\sin\dfrac{8\pi}7$ and $y=\cos\dfrac{2\pi}7+\cos\dfrac{4\pi}7+\cos\dfrac{8\pi}7$, then $\dfrac{5(x^2+y^2)}2$ is
\qn{56} ABCD is a trapezium. AB and CD are parallel and AD is perpendicular to them. Let AB = 4, CD = 3 and r be the radius of the circle which touches all the sides of trapezium them the value of [r] is (where [.] denotes greatest integer function)
\qn{57} Number of roots of the equation $|\sin\theta|+|\cos\theta|=\tan^2\theta+\cot^2\theta$
\qn{58} The sides a, b, c of $\triangle$ABC satisfy the relations $\mathrm a+\mathrm b-\mathrm c=2$ and $2\mathrm{ab}-\mathrm c^2=4$, if $x\in[0,1]$ then max value of $ax^2+bx+c$ is :
\qn{59} In parallelogram ABCD, P is any point inside parallelogram such that sum of area of $\triangle$APD and $\triangle$BPC is 3 sq. unit, then area of parallelogram ABCD is
\qn{60} The number of solutions of the equation $\tan x+\sec x=2\cos x$ lying in the interval $[0,2\pi]$ is :
\end{multicols}
\end{document}
